\chapter{Purpose and Scope}

This document is the implementation contract for the first real NavKit
inertial navigator. It is deliberately narrower than the general
\texttt{navigation\_reference} text. The goal is to define the first physically
meaningful, testable, and maintainable ECEF INS/GNSS path before writing the
mechanization code.

The primary reference convention is strapdown inertial navigation notation
\cite{groves2013,titterton2004}, with covariance and filtering support from
\cite{simon2006} and the existing NavKit navigation reference. Quaternion
notation and multiplicative attitude-error handling are informed by
\cite{sola2018}.

\begin{scopebox}
Version 1 targets a standard body-collocated IMU mechanization with GNSS
antenna lever-arm aiding:
\begin{itemize}
    \item one IMU;
    \item IMU frame collocated with the body/navigation state point;
    \item no lever arm between IMU and navigation state;
    \item configured lever arms for non-IMU aiding sensors, beginning with the
    GNSS antenna;
    \item ECEF/PCPF nominal mechanization;
    \item quaternion nominal attitude propagation;
    \item small-angle attitude error state;
    \item Kalman covariance/error-state propagation separate from the nominal quaternion state.
\end{itemize}
\end{scopebox}

\section{Primary Departure from Textbook DCM Mechanization}

Many textbook derivations present the nominal attitude mechanization directly
with direction cosine matrices. NavKit v1 will use a unit quaternion as the
nominal attitude state and derive the equivalent DCM from the quaternion when
rotating vectors. This is an implementation choice, not a different physical
model.

The error-state filter still carries a three-component small-angle attitude
perturbation. During correction, that small-angle vector is converted into a
small quaternion correction and applied multiplicatively to the nominal
quaternion. The quaternion is then normalized and the attitude error state is
reset.

\section{Explicit Non-Goals}

The following are important, but intentionally deferred:
\begin{itemize}
    \item multiple IMUs;
    \item arbitrary navigation reference point distinct from the IMU/body frame;
    \item IMU lever-arm correction;
    \item local-level and wander-azimuth mechanizations;
    \item PCI/ECI mechanization;
    \item estimating lever arms as calibration states;
    \item higher-order coning/sculling compensation beyond the first validated
    two-sample implementation;
    \item runtime-polymorphic mechanization selection.
\end{itemize}

\begin{implbox}
This document gates implementation. The first production code should implement
this scope before reintroducing the more general origin-frame and lever-arm
formulation from the broader navigation reference.
\end{implbox}
